\section{Solución Propuesta}
\label{sec:solucion}

\subsection{Objetivos}
\label{sec:objetivos}

\subsubsection{Objetivo general}

Diseñar, implementar y validar \STELLAR{} (\emph{Spectral Type Estimation via
Language Learning and Astronomical Reasoning}), un sistema de clasificación
espectral estelar híbrido que combine una capa de \emph{Hard Computing} (HC)
determinista con el modelo de lenguaje especializado \AstroSage{} (\emph{Soft
Computing}, SC), de modo que el anclaje determinista corrija el sesgo \emph{a
priori} del modelo de lenguaje y se obtenga una clasificación MK precisa,
reproducible y acompañada de una justificación astrofísica en lenguaje natural,
operando sobre \num{498} estrellas de Gaia DR3.

\subsubsection{Objetivos específicos}

\begin{enumerate}
  \item \textbf{Capa HC determinista.} Implementar los cuatro agentes
  deterministas del pipeline HC-2.0 (\texttt{AstrometryAgent},
  \texttt{ContinuumAgent}, \texttt{LineAgent} y \texttt{BinaryDetectorAgent})
  que, a partir de la metadata tabular de Gaia DR3, deriven los parámetros
  físicos limpios (\MG{}, \teff{}, $[M/H]$, $[Fe/H]$, $[\alpha/Fe]$,
  $V_\mathrm{tan}$) y un conjunto de banderas lógicas
  (\emph{is\_reliable\_parallax}, \emph{is\_giant}, \emph{is\_metal\_poor},
  \emph{is\_binary\_candidate}, \emph{is\_high\_velocity},
  \emph{has\_emission}).

  \item \textbf{Contrato de interfaz HC$\rightarrow$SC.} Definir y generar un
  contrato JSON por estrella que sirva de entrada estructurada y predigerida al
  modelo de lenguaje, respetando estrictamente las convenciones del catálogo
  (signo de H$\alpha$ ESP-ELS, tratamiento de \texttt{NaN} y del artefacto
  $\mathrm{fe\_h}=0.0$).

  \item \textbf{Capa SC con anclaje.} Configurar la inferencia de \AstroSage{}
  para que, recibiendo el contrato HC, emita la clasificación MK (tipo, rango de
  subtipo y clase de luminosidad), el grupo poblacional, puntuaciones de
  confianza y una descripción técnica, de forma que el anclaje HC acote el
  espacio de respuestas y mitigue el prior hacia el tipo G.

  \item \textbf{Optimización iterativa del \emph{prompt}.} Iterar el diseño del
  \emph{system prompt} de la capa SC a lo largo de siete versiones (V1--V7),
  aislando el efecto de cada decisión (anclaje, penalización por binariedad,
  recuperación RAG, límites de extensión) sobre las métricas de validación.

  \item \textbf{Validación multinivel.} Evaluar el sistema de extremo a extremo
  contra referencias externas: clasificación MK gruesa frente a SIMBAD;
  parámetros físicos (\teff{}, \logg{}) frente al catálogo PASTEL; calibración
  de confianza; y calidad de las descripciones en lenguaje natural (ROUGE,
  BERTScore y coherencia), reportando las métricas con intervalos de confianza
  \emph{bootstrap}.

  \item \textbf{Reproducibilidad.} Documentar el corpus ($n=498$), los supuestos,
  los hiperparámetros (versiones de \emph{prompt}) y la infraestructura de
  cómputo, de modo que los resultados sean verificables y el sistema extensible
  a muestras mayores de Gaia DR3/DR4.
\end{enumerate}

\subsection{Marco Teórico y Metodología}
\label{sec:marco-teorico}

Esta sección establece, primero, los fundamentos astrofísicos sobre los que se
sostiene la clasificación (§\ref{sec:fund-mk}--§\ref{sec:fund-lineas}) y, segundo,
el marco metodológico que articula las capas HC y SC
(§\ref{sec:metodo-arquitectura}--§\ref{sec:metodo-validacion}). El detalle de
implementación y la formulación matemática de cada agente determinista se
desarrolla en la Sección~\ref{sec:modelos}.

\subsubsection{Clasificación espectral MK y temperatura efectiva}
\label{sec:fund-mk}

El sistema de Morgan--Keenan ordena las estrellas en la secuencia
O--B--A--F--G--K--M según la temperatura efectiva \teff{} de su fotosfera, de las
más calientes (O, $\gtrsim\SI{30000}{\kelvin}$) a las más frías
(M, $<\SI{3700}{\kelvin}$). Cada letra se subdivide en diez subtipos (0--9) que
refinan la escala de temperatura dentro de la clase. En \STELLAR{}, la asignación
de la letra es competencia de la capa HC, que aplica umbrales estrictos sobre
\teff{} (\Cref{tab:teff-letra}). El subtipo, en cambio, es tarea de la capa SC,
que opera dentro de la letra ya fijada. La incertidumbre de
\SIrange{100}{200}{\kelvin} propia de las estimaciones \teff{} de GSP-Phot motiva
que el subtipo se exprese como un \emph{rango} (p.\,ej.\ \texttt{"3-5"}) y no como
un entero único.

\begin{table}[H]
  \centering
  \caption{Umbrales deterministas \teff{}$\rightarrow$letra MK aplicados por la
  capa HC. La frontera B/A se fija en exactamente \SI{10000}{\kelvin}.}
  \label{tab:teff-letra}
  \begin{tabular}{lc}
    \toprule
    Rango \teff{} (K) & Letra asignada \\
    \midrule
    $\ge \num{30000}$        & O \\
    \numrange{10000}{29999}  & B \\
    \numrange{7500}{9999}    & A \\
    \numrange{6000}{7499}    & F \\
    \numrange{5200}{5999}    & G \\
    \numrange{3700}{5199}    & K \\
    $< \num{3700}$           & M \\
    \bottomrule
  \end{tabular}
\end{table}

La frontera B/A en \SI{10000}{\kelvin} es la zona de confusión dominante del
sistema. Para mitigarla, la versión V7 incorpora una corrección por gravedad: las
estrellas con \teff{} en \SIrange{9700}{10100}{\kelvin} y $\log g \ge 3.8$ se
asignan a~A (régimen enana/subgigante), regla que eleva la exactitud sobre
estrellas frontera de \SI{43.1}{\percent} (umbral \teff{} estricto) a
\SI{76.5}{\percent}.

\subsubsection{Clase de luminosidad y gravedad superficial}
\label{sec:fund-lum}

La clase de luminosidad (I a V) distingue el estado evolutivo a igual temperatura:
desde supergigantes (I) hasta enanas de secuencia principal (V). Uno de sus 
discriminantes físicos directos es la gravedad superficial \logg{}: las gigantes 
evolucionadas, de gran radio, tienen gravedades bajas, mientras que las enanas 
compactas exhiben gravedades altas. Por ello \STELLAR{} delega la clase de 
luminosidad enteramente a la capa SC, alimentándola con el valor crudo 
\verb|logg_gspphot|, que evita toda transformación en HC. La \Cref{tab:logg-lum} 
resume el mapeo empleado.

\begin{table}[H]
  \centering
  \caption{Mapeo \logg{}$\rightarrow$clase de luminosidad sobre valores GSP-Phot
  del corpus. La zona subgigante (IV) es la de mayor ambigüedad.}
  \label{tab:logg-lum}
  \begin{tabular}{llc}
    \toprule
    Clase & Denominación & Rango \logg{} (dex) \\
    \midrule
    I   & Supergigante          & $< 1.5$ \\
    II  & Gigante brillante     & \numrange{1.5}{2.5} \\
    III & Gigante               & \numrange{2.5}{3.5} \\
    IV  & Subgigante            & \numrange{3.5}{4.0} \\
    V   & Enana (sec.\ principal) & $> 4.0$ \\
    \bottomrule
  \end{tabular}
\end{table}

Dado que \logg{} llega sin procesar, la capa SC debe contemplar sesgos
sistemáticos conocidos de GSP-Phot: la sobreestimación de \logg{} en gigantes
frías K/M (tratar el valor como cota inferior) y su menor fiabilidad en estrellas
calientes B/A (donde conviene contrastar con la magnitud absoluta \MG{}).

\subsubsection{Poblaciones galácticas: química y cinemática}
\label{sec:fund-pob}

La Vía Láctea alberga poblaciones estelares con historias de formación
distinguibles por su química y cinemática. \STELLAR{} discrimina tres grupos:

\begin{itemize}
  \item \textbf{Disco Fino:} poblaciones jóvenes a intermedias, metalicidad
  cercana a la solar ($[M/H] \ge -0.2$), baja velocidad espacial.
  \item \textbf{Disco Grueso:} estrellas viejas, cinemáticamente calientes,
  metalicidad subsolar ($-1.0 \le [M/H] < -0.2$) con realce de elementos
  $\alpha$ ($[\alpha/Fe] \ge 0.2$), reliquia de una época de formación estelar
  rápida y temprana.
  \item \textbf{Halo:} estrellas muy viejas y pobres en metales
  ($[M/H] < -1.0$), o bien con alta velocidad tangencial
  ($V_\mathrm{tan} > \SI{200}{\kilo\meter\per\second}$).
\end{itemize}

La población se asigna en HC mediante una \emph{cadena de prioridad}: el primer
criterio que se cumple decide y los posteriores no se evalúan. La cinemática
($V_\mathrm{tan}$) tiene prioridad sobre la química porque es un indicador más
directo de pertenencia al Halo, le sigue la regla de realce $\alpha$ para el
Disco Grueso y, finalmente, los umbrales de metalicidad. La química se resuelve
priorizando el valor espectroscópico $[Fe/M]$ (\verb|fem_gspspec|, $\sim$44\%
de cobertura) sobre el fotométrico $[M/H]$ (\verb|mh_gspphot|), el valor
$\mathrm{fe\_h}=0.0$ exacto se trata como artefacto nulo y se recurre al
fotométrico.

\subsubsection{Magnitudes derivadas, extinción y diagrama HR}
\label{sec:fund-hr}

A partir de la astrometría de Gaia se derivan dos cantidades físicas clave. La
magnitud absoluta en banda~G se obtiene del módulo de distancia corregido por
extinción interestelar,
\begin{equation}
  M_G = G + 5 + 5\log_{10}\!\left(\frac{\varpi}{1000}\right) - A_G,
  \qquad A_G = 2.74\,E(BP\!-\!RP),
  \label{eq:absmag}
\end{equation}
donde $\varpi$ es el paralaje en mas y $A_G$ la extinción en banda~G. La velocidad
tangencial se obtiene de los movimientos propios,
\begin{equation}
  V_\mathrm{tan} = \frac{4.74\,\sqrt{\mu_{\alpha*}^2 + \mu_\delta^2}}{\varpi}
  \quad [\si{\kilo\meter\per\second}].
  \label{eq:vtan}
\end{equation}
Ambas magnitudes sólo son físicamente válidas cuando el paralaje es confiable
(SNR $=\varpi/\sigma_\varpi > 5$); en caso contrario se suprimen del contrato. La
combinación de \MG{} y \teff{} sitúa a la estrella en el diagrama
Hertzsprung--Russell, sobre el que descansa la distinción gigante/enana.

\subsubsection{Diagnósticos espectrales de línea}
\label{sec:fund-lineas}

Dos diagnósticos espectrales complementan la fotometría. El \textbf{triplete de
Ca\,\textsc{ii}} (849.8, 854.2, \SI{866.2}{\nano\meter}), presente en la ventana
RVS de Gaia, es altamente sensible a la gravedad superficial y la metalicidad. La
capa HC ajusta perfiles de Voigt sobre estas líneas para extraer ancho equivalente
y anchura. La línea de \textbf{H$\alpha$} (\SI{656.3}{\nano\meter}) es un trazador
de actividad cromosférica y emisión circumestelar, al quedar fuera de la ventana
RVS y ser inalcanzable a la resolución BP/RP, se toma directamente del catálogo
ESP-ELS (\verb|ew_espels_halpha|). Aquí rige una convención de signo crítica:
\begin{equation}
  \texttt{ew\_espels\_halpha} > 0 \Rightarrow \text{absorción}, \qquad
  \texttt{ew\_espels\_halpha} < 0 \Rightarrow \text{emisión},
  \label{eq:halpha}
\end{equation}
de modo que \emph{has\_emission} se activa cuando el ancho equivalente es negativo
y la bandera de calidad es 0. El respeto estricto de esta convención evita errores
sistemáticos en la identificación de estrellas activas, Be o simbióticas.

\subsubsection{Hard Computing, Soft Computing y modelos de lenguaje}
\label{sec:metodo-paradigmas}

El trabajo se enmarca en la complementariedad de dos paradigmas de cómputo. El
\emph{Hard Computing} agrupa los métodos deterministas, exactos y verificables
(física clásica, álgebra, estadística) idóneos para extraer parámetros y
aplicar reglas físicas duras. El \emph{Soft Computing} agrupa los métodos
tolerantes a la imprecisión y orientados al razonamiento aproximado, en este
trabajo lo encarna un modelo de lenguaje a gran escala. El uso de \AstroSage{}
(un LLM de \num{8}\,mil millones de parámetros afinado sobre literatura
astronómica) aporta la capacidad de integrar evidencia heterogénea y de
articular el razonamiento en lenguaje natural, a costa de los riesgos propios de
los LLM (alucinación, sensibilidad al formato y sesgo \emph{a priori}). La
hipótesis metodológica de \STELLAR{} es que estos riesgos pueden acotarse
anclando la decisión del LLM con la salida verificable del HC.

\paragraph{El modelo AstroSage-8B.} \AstroSage{} (formalmente
AstroSage-Llama-3.1-8B) es un asistente de inteligencia artificial de código
abierto especializado en astronomía, astrofísica, cosmología e instrumentación
astronómica, construido sobre la arquitectura de Meta-Llama-3.1-8B
\citep{astrosage}. Su desarrollo se divide en construcción, entrenamiento y
validación.

\emph{Construcción.} El modelo conserva la arquitectura base y el tamaño de su
modelo fundacional (\num{8}\,mil millones de parámetros) y mantiene el
tokenizador original de Llama-3.1, sin añadir \emph{tokens} de vocabulario
astronómico específicos. Su entrenamiento se ejecutó sobre el superordenador de
clase exaescala \emph{Frontier} del \emph{Oak Ridge Leadership Computing
Facility} (OLCF), empleando cientos de nodos con GPU AMD MI250X.

\emph{Entrenamiento.} Se realizó en tres etapas consecutivas:
\begin{itemize}
  \item \textbf{Preentrenamiento continuo (\emph{Continued Pre-training}, CPT).}
  Infunde conocimiento del dominio a partir de un corpus masivo de
  $\sim$\num{250000} \emph{preprints} de arXiv (categorías \texttt{astro-ph} y
  \texttt{gr-qc}, 2007--enero 2024), $\sim$\num{30000} artículos de Wikipedia de
  astronomía y $\sim$\num{800} libros de texto, totalizando \num{3.3}\,mil
  millones de \emph{tokens} (\SI{19.9}{\giga\byte} de texto). El texto se extrajo
  de los PDF con herramientas OCR (Nougat) y se depuró mediante un filtrado por
  \emph{perplejidad} que eliminó el \SI{2}{\percent} de los párrafos
  estadísticamente anómalos.
  \item \textbf{Ajuste fino supervisado (\emph{Supervised Fine-Tuning}, SFT).}
  Enseña al modelo a seguir instrucciones en formato pregunta--respuesta sobre
  \num{2}\,mil millones de \emph{tokens} (\SI{9.8}{\giga\byte}). El núcleo del
  conjunto son \num{8.8}\,millones de pares pregunta--respuesta sintéticos,
  generados a partir de artículos científicos y filtrados rigurosamente por otra
  IA según corrección, relevancia y calidad; se complementó con instrucciones de
  dominio general (Infinity-Instruct-7M) para refinar la habilidad
  conversacional.
  \item \textbf{Fusión de modelos (\emph{Model Merging}).} Para recuperar el
  pulido conversacional sin sacrificar el conocimiento astronómico, se promediaron
  los pesos mediante la técnica DARE-TIES (librería \texttt{mergekit}),
  combinando \SI{75}{\percent} del modelo AstroSage afinado con
  \SI{25}{\percent} de Meta-Llama-3.1-8B-Instruct.
\end{itemize}

\emph{Validación.} El rendimiento se verificó por tres vías:
\begin{itemize}
  \item \textbf{\emph{Benchmark} AstroMLab-1.} Batería de \num{4425} preguntas de
  opción múltiple creadas por expertos a partir de artículos de \emph{Annual
  Review of Astronomy and Astrophysics}, deliberadamente excluidos del corpus de
  entrenamiento para evitar memorización. AstroSage-8B alcanzó una exactitud del
  \SI{80.9}{\percent}, superando en 8 puntos a su base Llama-3.1-8B y empatando con
  modelos corporativos como GPT-4o (\SI{80.4}{\percent}) a una fracción del coste.
  \item \textbf{Métricas generales.} En \emph{benchmarks} de conocimiento general y
  matemáticas (MATH, BBH, MMLU-PRO, GPQA) se comprobó que la fusión de modelos no
  indujo \emph{olvido catastrófico}, preservando su capacidad de razonamiento
  lógico fuera de la astronomía.
  \item \textbf{Evaluación humana a ciegas.} Expertos del área compararon respuestas
  de AstroSage-8B frente a Meta-Llama-3.1-8B-Instruct sin conocer su autoría; en el
  \SI{73}{\percent} de los casos prefirieron las de AstroSage por su mayor exactitud
  y calidad, especialmente en cosmología.
\end{itemize}

Estas propiedades ---conocimiento astronómico profundo y razonamiento preservado---
hacen de \AstroSage{} un candidato idóneo para la capa SC de \STELLAR{}, en la que
opera no como clasificador autónomo sino anclado por la evidencia verificable del HC.

\subsubsection{Arquitectura híbrida y flujo de trabajo}
\label{sec:metodo-arquitectura}

\STELLAR{} se organiza en dos capas encadenadas por un contrato JSON
(\Cref{fig:pipeline}):

\begin{enumerate}
  \item \textbf{Capa HC (pipeline HC-2.0).} Cuatro agentes deterministas
  (\texttt{AstrometryAgent}, \texttt{ContinuumAgent}, \texttt{LineAgent} y
  \texttt{BinaryDetectorAgent}) procesan la metadata tabular y los espectros
  de cada estrella. Producen un \emph{contrato} por estrella que contiene un
  vector físico ($M_G$, \teff{}, $[M/H]$, $[Fe/H]$, $[\alpha/Fe]$, \logg{},
  $V_\mathrm{tan}$, $A_G$), un conjunto de banderas lógicas booleanas, los
  diagnósticos de binariedad y espectrales, y un \emph{quality\_score} global.
  \item \textbf{Capa SC (\AstroSage{}).} Recibe el contrato como entrada
  predigerida y emite un JSON estructurado con la clasificación (tipo, rango de
  subtipo, clase de luminosidad, grupo poblacional), puntuaciones de confianza y
  una justificación técnica en lenguaje natural.
\end{enumerate}

\begin{figure}[H]
  \centering
  \resizebox{\textwidth}{!}{%
  \begin{tikzpicture}[
    node distance=1.0cm and 0.5cm,
    % Datos de entrada / base de conocimiento -> ELIPSE GRIS
    data/.style={ellipse, draw=gray!80, fill=gray!5, text width=2.2cm,
      text centered, font=\sffamily\scriptsize, minimum height=1.1cm, thick},
    % Agentes Hard Computing (deterministas) -> AZUL
    hard/.style={rectangle, draw=blue!50!black, fill=blue!10, text width=2.7cm,
      text centered, rounded corners=4pt, minimum height=1.5cm,
      font=\sffamily\scriptsize\bfseries, drop shadow},
    % Contrato JSON (handoff HC -> SC) -> MORADO
    contract/.style={rectangle, draw=violet!60!black, fill=violet!10, text width=9cm,
      text centered, rounded corners=4pt, minimum height=1.1cm,
      font=\sffamily\scriptsize, drop shadow},
    % Capa Soft Computing (AstroSage-8B) -> VERDE
    soft/.style={rectangle, draw=green!40!black, fill=green!10, text width=9cm,
      text centered, rounded corners=4pt, minimum height=1.4cm,
      font=\sffamily\small\bfseries, drop shadow},
    % Resultado final -> AMARILLO
    result/.style={rectangle, draw=yellow!40!black, fill=yellow!15, text width=8cm,
      text centered, font=\sffamily\footnotesize\itshape, drop shadow, rounded corners,
      minimum height=1.1cm},
    arrow/.style={-{Stealth[scale=1.0]}, thick, color=gray!70!black},
    layerbox/.style={draw=gray!30, dashed, fill=gray!5, rounded corners=10pt, inner sep=10pt}
  ]

    % --- 1. AGENTES HC (fila central de referencia) ---
    \node (cont)   [hard] {ContinuumAgent\\{\sffamily\tiny\mdseries spline + $\sigma$-clipping}};
    \node (binary) [hard, left=of cont] {BinaryDetectorAgent\\{\sffamily\tiny\mdseries RUWE · RV · NSS}};
    \node (astro)  [hard, left=of binary] {AstrometryAgent\\{\sffamily\tiny\mdseries $M_G$ · $V_\mathrm{tan}$ · flags}};
    \node (line)   [hard, right=of cont] {LineAgent\\{\sffamily\tiny\mdseries Voigt · Ca\,\textsc{ii} · EW}};

    % --- 2. ENTRADAS ---
    \node (tab)  [data, above=1.5cm of binary] {Tabular CSV\\Gaia DR3\\{\tiny(astrometría, fotometría, química)}};
    \node (rvs)  [data, above=1.5cm of cont] {Espectros RVS\\{\tiny846--870 nm}};
    \node (bprp) [data, above=1.5cm of line] {Espectros BP/RP\\{\tiny336--1020 nm}};

    % --- 3. CONTRATO (handoff) — centrado bajo el punto medio de la fila HC ---
    \coordinate (hcmid) at ($(binary.south)!0.5!(cont.south)$);
    \node (contract) [contract, below=1.7cm of hcmid]
      {\textbf{Contrato JSON (HC-2.0)}\\
       physical\_vector + logical\_flags + spectral\_summary + \texttt{hc\_anchor} + quality\_score};

    % --- 4. CAPA SC (AstroSage-8B) ---
    \node (sc) [soft, below=1.1cm of contract] {Capa SC: \AstroSage{}\\{\sffamily\tiny\mdseries inferencia + subtipo · luminosidad · confianzas · descripción}};

    % --- Base de conocimiento (RAG) ---
    \node (kb) [data, right=1.6cm of sc] {Base de conocimiento\\{\tiny(guías MK, luminosidad, poblaciones, flags)}};

    % --- 5. SALIDA ---
    \node (out) [result, below=1.1cm of sc]
      {JSON final: clasificación MK (tipo · subtipo · luminosidad · población) + confianzas + razonamiento};

    % --- Agrupaciones visuales por capa ---
    \begin{scope}[on background layer]
      \node [layerbox, fit=(astro)(binary)(cont)(line),
        label={[gray!70, font=\sffamily\scriptsize\bfseries]north:CAPA HARD COMPUTING (HC-2.0)}] {};
    \end{scope}

    % --- Conexiones: entradas -> agentes ---
    \draw [arrow] (tab) -- (astro);
    \draw [arrow] (tab) -- (binary);
    \draw [arrow] (rvs) -- (cont);
    \draw [arrow] (bprp) -- (cont);
    \draw [arrow] (cont) -- node[above, font=\tiny, color=gray] {flujo norm.} (line);

    % --- Agentes -> contrato ---
    \draw [arrow] (astro.south)  to[out=270, in=170] (contract.west);
    \draw [arrow] (binary.south) -- (contract.north -| binary.south);
    \draw [arrow] (cont.south)   -- (contract.north -| cont.south);
    \draw [arrow] (line.south)   to[out=270, in=10] (contract.east);

    % --- Contrato -> SC -> salida ---
    \draw [arrow] (contract) -- node[right, font=\tiny\itshape, color=gray] {anclaje \texttt{hc\_anchor}} (sc);
    \draw [arrow] (kb) -- node[above, font=\tiny\itshape, color=gray] {RAG top-3} (sc);
    \draw [arrow] (sc) -- (out);

  \end{tikzpicture}%
  }
  \caption{Arquitectura de dos capas de \STELLAR{}. Los cuatro agentes
  deterministas de la capa HC procesan la metadata tabular y los espectros de
  Gaia DR3 y consolidan un contrato JSON por estrella; este, a través del campo
  \texttt{hc\_anchor}, ancla la inferencia de la capa SC (\AstroSage{}), que
  ---apoyada por recuperación aumentada (RAG) sobre la base de conocimiento---
  emite la clasificación MK final y su justificación en lenguaje natural.}
  \label{fig:pipeline}
\end{figure}

El contrato HC, además del schema completo, transporta un campo \verb|hc_anchor|
que comunica a la capa SC las decisiones ya tomadas de forma determinista.

\subsubsection{Estrategia de anclaje (\texttt{hc\_anchor})}
\label{sec:metodo-anclaje}

El mecanismo central que corrige el sesgo del LLM es el \emph{anclaje}: cuando el
contrato incluye un \verb|hc_anchor|, la \textbf{letra espectral} y el
\textbf{grupo poblacional} ya han sido fijados por las reglas deterministas
(\Cref{tab:teff-letra} y la cadena de prioridad de §\ref{sec:fund-pob}) y la capa
SC tiene prohibido recalcularlos o sobrescribirlos. La tarea del LLM se restringe
entonces a (i)~asignar el subtipo dentro de la letra fijada, (ii)~asignar la clase
de luminosidad a partir de \logg{}, (iii)~justificar la población, y (iv)~emitir
confianzas y descripción. El LLM puede expresar desacuerdo únicamente mediante
confianzas más bajas, nunca cambiando la etiqueta anclada. Este diseño impide que
el prior hacia el tipo~G distorsione la clasificación y convierte al LLM en un
motor de refinamiento y explicación sobre un esqueleto verificable.

\subsubsection{Optimización iterativa del \emph{system prompt}}
\label{sec:metodo-prompt}

La capa SC se desarrolla de forma incremental a lo largo de siete versiones del
\emph{system prompt} (V1--V7). Cada versión aísla el efecto de una decisión de
diseño sobre las métricas de validación: introducción del anclaje, corrección de
la frontera A/B por \logg{}, penalización de confianza ante binariedad, límites de
extensión de la descripción y la incorporación de recuperación aumentada
(RAG) sobre una base de conocimiento astronómico. El detalle de cada versión y su
impacto se presenta en la Sección~\ref{sec:resultados}.

\subsubsection{Protocolo de validación multinivel}
\label{sec:metodo-validacion}

La credibilidad del sistema se evalúa de extremo a extremo contra referencias
externas independientes:

\begin{itemize}
  \item \textbf{Nivel 1 --- MK gruesa (vs.\ SIMBAD).} Se contrasta la letra
  espectral contra el tipo \verb|sp_type| de SIMBAD (cobertura 491/498). Métricas:
  exactitud, Cohen $\kappa$, F1 macro y ponderado, distancia MK media y exactitud
  \emph{near-miss} ($d\le1$), con intervalos de confianza \emph{bootstrap}.
  \item \textbf{Nivel 2 --- Parámetros físicos (vs.\ PASTEL).} Se compara
  \teff{} y \logg{} inferidos contra el catálogo PASTEL de alta resolución
  (cobertura 340 y 194 estrellas respectivamente). El campo
  \verb|n_pastel_measurements| pondera la solidez de cada referencia.
  \item \textbf{Nivel 3 --- Calibración de confianza.} Se examina la correlación
  entre las confianzas emitidas y la corrección efectiva de la predicción.
  \item \textbf{Nivel 4 --- Calidad de la descripción.} Se evalúan las
  descripciones en lenguaje natural mediante ROUGE y BERTScore frente a
  referencias, coherencia interna y exactitud de subtipo.
  \item \textbf{Nivel 5 --- Impacto del RAG.} Se mide la tasa de acierto de la
  recuperación y su efecto sobre la calidad de las descripciones.
\end{itemize}

Esta batería multinivel cierra el lazo de validación científica, distinguiendo
explícitamente la calidad de la \emph{clasificación} (niveles 1--2) de la calidad
de la \emph{explicación} (niveles 4--5).

% TODO: figura del pipeline conceptual HC->contrato->SC (figuras/pipeline.pdf)
% TODO: citas formales — Wang & Chen (2019, A_G), Lindegren et al. (2021, RUWE),
%       Soubiran et al. (PASTEL), Gaia DR3, AstroSage-8B.
